\section{Methodology}
\label{ch:methodology}

This section describes the experimental platform in the detail needed to reproduce the study, and it does so in a particular order: it first establishes the task and the models, then the samplers and the metrics, and only then the diagnostic experiments, because the diagnostics were designed after the main experiments and in response to what those experiments revealed. Where a choice was made, the reasoning behind it is given, because in a diagnostic study the design of each measurement is part of the argument. The section sets out the masked-recovery task and why it was chosen, the five energy functions, the two samplers and the factorial space of configurations, the evaluation metrics and the justification for each, the GFlowNet training setup, the constrained-generation extension, and finally the five diagnostic experiments, each introduced together with the question it was built to answer.

\subsection{Task: Masked Token Recovery}
\label{sec:meth-task}

A sampler can only be studied if there is a way to tell whether it is being guided toward good sequences or is merely wandering. An open-ended generation task provides no such handle, because there is no reference against which to judge the output: a fluent but arbitrary continuation is neither clearly a success nor clearly a failure of the sampler, and one cannot separate the sampler's contribution from the difficulty of the underlying generation. The study therefore adopts \textbf{masked token recovery}, also called infilling, as its diagnostic bench. A sequence is drawn from the corpus, one or more of its tokens are replaced with random tokens, and the sampler is asked to restore a coherent sequence by resampling only the corrupted positions while the rest of the sequence is held fixed. Because the original sequence is known, the recovery can be scored; because the corruption and the surrounding context are controlled, the behaviour of the sampler can be attributed to the method rather than to the open-endedness of the task. This is the standard reason infilling is used as a diagnostic for controllable generation, and it is the reason it is used here.

An important subtlety shapes the choice of evaluation metric and is worth stating at the outset. The primary metric, defined in Section~\ref{sec:meth-metrics}, compares the next-token distribution that the recovered fill induces against the distribution induced by the ground-truth fill, so it measures how well the recovered token fits its context relative to a reference, and a coherent synonym of the original token, such as \emph{dorm} for \emph{roommate}, scores well even though it does not match. This is a measure of relative-to-reference contextual fit, not of absolute sequence likelihood used as a proxy for quality. It is therefore not in tension with the likelihood-trap finding of Section~\ref{sec:results-trap}, which concerns the absolute sequence likelihood as a maximization objective in open-ended generation, a different quantity serving a different purpose. The full justification for preferring this metric over exact-match accuracy is given once, with the metric definition, in Section~\ref{sec:meth-metrics}.

Two corpora are used, and it matters which experiment ran on which. The masked-recovery sequences of the main sampling grid are drawn from the \textbf{WikiText-2} validation set \citep{merity2017pointer}, a standard held-out language-modelling corpus of encyclopedic prose, chosen because its held-out status keeps the recovery task independent of any fine-tuning data and its sentences give a masked position a locally checkable correct recovery. The base GPT-2 Large is instead supervised fine-tuned on \textbf{ROCStories} \citep{mostafazadeh2016corpus}, a collection of short five-sentence everyday narratives and the corpus of the amortized-inference reference implementation of \citet{hu2024amortizing}, so the GFlowNet half of the study trains and generates on ROCStories and operates on the same data distribution as the work it builds on. The main sampling grid therefore evaluates recovery \emph{out of domain} relative to the fine-tuning corpus, while the diagnostic experiments are run in domain on ROCStories. Table~\ref{tab:corpus-map} maps each experiment family to its corpus.

\begin{table}[htbp]
    \centering
    \small
    \begin{tabular}{ll}
        \toprule
        Experiment family & Corpus \\
        \midrule
        Masked-recovery grid (145 configurations) & WikiText-2 validation \\
        Gradient-free baselines, continuation, external judge & WikiText-2 validation \\
        Constrained generation & WikiText-2 validation \\
        Diffusion positive control (SEDD) & WikiText-2 validation \\
        GPT-2 Large fine-tuning; GFlowNet training and variants & ROCStories \\
        Diagnostics: linearization, likelihood trap, anisotropy, trajectory & ROCStories \\
        \bottomrule
    \end{tabular}
    \caption[Corpus used by each experiment family.]{Corpus used by each experiment family. The base model is fine-tuned on ROCStories; the main sampling grid recovers out-of-domain WikiText-2 sequences, while the diagnostics run in domain on ROCStories.}
    \label{tab:corpus-map}
\end{table}

The domain shift is not innocent, so it is addressed directly. If the gradient null were an artifact of evaluating out-of-domain text it should weaken or vanish in domain, and it does not: the in-domain ROCStories diagnostics of Section~\ref{sec:results-linradius} and Section~\ref{sec:results-trap} reproduce both the gradient null, with the linearization correlation below $0.06$ throughout, and the likelihood trap. The central findings are therefore properties of the frozen likelihood rather than of the domain shift.

\subsection{Energy Functions}
\label{sec:meth-models}

The study evaluates five energy functions, chosen to let the central findings be tested for robustness across training regimes and across architectures.

The reference energy is a \textbf{GPT-2 Large} model \citep{radford2019language}, of $774$ million parameters and embedding dimension $D = 1280$, supervised fine-tuned on ROCStories. This same fine-tuned model serves as the base for the amortized experiments, so that the Langevin and GFlowNet halves of the study share weights and are not confounded by any difference in the underlying network, an identity that is essential to the unifying experiment of Section~\ref{sec:results-gfn}. Three \textbf{GFlowNet-tuned variants} of this base model are included, produced as described in Section~\ref{sec:meth-gfn}: two trained with the unnormalized-length reward setting, at $500$ and $2000$ optimization steps respectively, so that the effect of longer training can be observed, and one trained with the length-normalized setting at $500$ steps. These tuned variants are included to test whether amortization repairs the energy landscape: the GFlowNet objective demonstrably moves the energy, as Section~\ref{sec:results-gfn} quantifies, and running the same samplers on the tuned models asks whether that movement made the landscape navigable by gradient or merely relocated it. Finally, a \textbf{Llama-3 8B} model \citep{dubey2024llama}, with embedding dimension $D = 4096$, serves as a cross-architecture control. Because Llama-3 uses a different tokenizer and a markedly different embedding geometry, and therefore a different step-size scale, it is never placed on a shared numerical axis with the GPT-2 family; it is used to establish that the qualitative findings, such as the gradient null of Section~\ref{sec:results-fallacy} and the Metropolis--Hastings breakdown, hold across architectures rather than being artifacts of one model. The four GPT-2-based models, by contrast, share the tokenizer, the corruption seeds, and the step-size schedule, and their numbers are therefore directly comparable without qualification.

\subsection{Samplers and Configurations}
\label{sec:meth-configs}

The two samplers of Section~\ref{sec:bg-samplers}, the Discrete Langevin Sampler (DLS) and the Continuous Langevin Sampler (CLS), are each run under a factorial combination of design choices. The two are studied together, rather than either alone, because they represent the two opposite strategies for confronting the discreteness of text and therefore isolate different failure modes: the DLS, following \citet{grathwohl2021oops} and \citet{zhang2022langevin}, keeps the state in token space and uses the gradient only to build a proposal, which tests whether the gradient is an informative ranking signal; the CLS, following \citet{kumar2022gradient}, relaxes into continuous space and runs genuine Langevin dynamics, which tests whether continuous Markov chain Monte Carlo is compatible with a projected discrete energy. A finding that holds for both is a finding about the energy rather than about one sampler's idiosyncrasies, which is why both are needed.

The continuous sampler's update as implemented deserves to be stated precisely, because Section~\ref{sec:results-traj} quantifies how far its state drifts from the token manifold and the update is what bounds that drift. At step $t$ with step size $\epsilon_t$, the sampler first forms an interim continuous state by the Langevin drift, $\vs_{\text{interim}} = \vs_t + \tfrac{\epsilon_t}{2}\,\grad_{\vs}\log p$, then projects it to the nearest token embedding, $\mathrm{proj}_V(\vs_{\text{interim}})$, and centres the proposal not at the interim state but at the interpolated mean of the two,
\begin{equation}
    \vs_{t+1} \sim \mathcal{N}\!\left(\tfrac{1}{2}\bigl(\vs_{\text{interim}} + \mathrm{proj}_V(\vs_{\text{interim}})\bigr),\; \epsilon_t \bm{I}\right),
    \label{eq:meth-cls-update}
\end{equation}
after which the Metropolis--Hastings correction of Section~\ref{sec:bg-mh} accepts or rejects it, with the reverse proposal density evaluated from the drift at the proposed point. Averaging the interim state with its projection pulls the proposal halfway back toward the token manifold at every step, which bounds the off-manifold drift a pure continuous update would accumulate. This interpolation is an implementation choice adapted from the continuous-relaxation samplers of \citet{kumar2022gradient}, not a property of Langevin dynamics itself; it is load-bearing for the off-manifold distances of Section~\ref{sec:results-traj}, since even with this pull toward the manifold the state settles far from any token.
% SOURCE: core/cls.py _step: interim = s + 0.5*eps*grad_s; interim_proj = emb[proj_V(interim)];
% m_s = 0.5*(interim + interim_proj); s_prop ~ N(m_s, eps I). Documented here 2026-07-26

One asymmetry in how the correction was applied across the compared arms has to be declared, because it falls on exactly the contrast this thesis rests on. In the main grid the reverse-proposal term of Equation~\eqref{eq:bg-mh} is computed exactly for the \emph{policy} arm and set to zero for the two random-direction arms, on the reasoning that a random-direction proposal is a symmetric random walk whose proposal ratio cancels. That reasoning is not sound. For the continuous sampler the proposal mean of Equation~\eqref{eq:meth-cls-update} carries both a drift and a nearest-neighbour projection, neither of which is symmetric, whatever the direction; for the discrete sampler the distance term is symmetric but the softmax normalizers at the current and proposed tokens are not. The consequence is that in every correction-enabled comparison the policy arm is charged a reversibility penalty from which its comparators are exempted, and Section~\ref{sec:results-mh} measures that penalty reaching $-1325$ nats in the continuous case. This works against the policy arm, so it cannot manufacture the null reported in Section~\ref{sec:results-fallacy}, but it could inflate the cases in which the gradient appears \emph{worse} than a random direction.

The samplers therefore carry a corrected mode, enabled by \texttt{-{}-mh\_exact\_all\_arms}, in which the reverse term is computed exactly for every arm. Because a random direction is drawn independently of the current state, holding that draw fixed across the forward and reverse evaluation of a single step yields a kernel $q_{\vg}$ that is exactly computable and $\pi$-invariant, and a mixture of $\pi$-invariant kernels is itself $\pi$-invariant, so the corrected mode is a valid sampler rather than an approximation. Section~\ref{sec:results-sweep} reports the arms rerun under it. The archived grid is left as it was run, and the legacy treatment remains the default, so that every number reported from it remains reproducible.
% SOURCE: core/dls.py and core/cls.py, the mh_exact_all_arms branch; the legacy branch is retained
% as the default and is what results/grid/{gpt2_v2,llama,gfn} were produced with.
% (author issue list item 9); previously undocumented and load-bearing for Section 5.4.

The design axes are as follows. The \textbf{proposal method} is one of three. The \emph{policy} method uses the true gradient of the energy. The \emph{grad-norm-preserved random} method replaces the gradient with a random vector rescaled to the exact norm of the true gradient, so that it differs from the policy only in direction. The \emph{random} method uses a random vector of random magnitude. The comparison among these three is the central ablation of the thesis, isolating respectively the full gradient, the effect of its direction alone, and a pure baseline, and the reasoning behind this three-way design is developed where it is used, in Section~\ref{sec:results-fallacy}. The \textbf{Metropolis--Hastings correction} is either enabled or disabled, so that the effect of theoretical correctness can be observed directly. \textbf{Gradient normalization}, which rescales the gradient to unit norm for numerical stability, is either enabled or disabled; this choice interacts with the proposal-method ablation in a way that turns out to be critical and is explained in Section~\ref{sec:results-fallacy}. The \textbf{annealing schedule} runs for either $50$ or $100$ steps, with the step size decreasing linearly; for the GPT-2 family the schedule runs from $\epsilon_{\text{start}} = 10.5$ to $\epsilon_{\text{end}} = 0.1$, a scale established by the calibration described in Section~\ref{sec:results-stepsize}. An \emph{oracle} variant, in which the ideal step size at each iteration is selected using knowledge of the ground-truth token, is included as an upper-bound reference; its purpose is to rule out the possibility that the samplers fail merely for want of a well-tuned schedule, since if even an oracle-tuned schedule does not let the gradient outperform a random direction, the failure cannot be one of calibration. The oracle sweep is run independently for each proposal method, with a separate per-method oracle run rather than a single schedule shared across methods, so the schedule selection cannot advantage the policy gradient over the random comparators. In total, across the five energy functions, the study comprises $145$ sampling configurations, each evaluated on $200$ sequences; the arithmetic of that count is given in Appendix~\ref{app:grid}.

\subsection{Evaluation Metrics}
\label{sec:meth-metrics}

Three quantities are recorded, and the choice among them is itself informative.

The first is the \textbf{Euclidean distance} in embedding space between the recovered token and the ground-truth token. This is the most obvious measure of recovery, and one might expect it to be the natural criterion. It proves to be unreliable, for a reason established quantitatively in Section~\ref{sec:results-aniso}: because language-model embeddings are anisotropic, occupying a narrow cone in the representation space \citep{ethayarajh2019contextual}, Euclidean distance compresses the distinction between good and bad tokens, so that a token can be geometrically close to the target while being grammatically inadmissible, and conversely a perfectly good synonym can lie far away. It is reported for completeness, and studied as an object in its own right, but it is not relied upon as the primary criterion.

The primary criterion is the \textbf{KL divergence} \citep{kullback1951information} between the model's predictive distribution at the recovered position and a reference distribution, which measures how well the recovered token fits its context, that is, its \emph{contextual coherence}. Concretely, let $\mathcal{M}$ be the set of masked positions in a sequence of length $L$, and let $\mathcal{M}' = \{\, m \in \mathcal{M} : m < L-1 \,\}$ be those masked positions that have a right neighbour, the only ones at which the metric is defined. For a position $m \in \mathcal{M}'$, write $p^{(m)}_{\text{ref}}$ for the model's next-token distribution at $m$ when every masked position holds its \emph{ground-truth} token, and $p^{(m)}_{\text{pred}}$ for the same distribution when the masked positions instead hold the \emph{recovered} fill. The metric is the mean divergence over these positions,
\begin{equation}
    \overline{\KL} = \frac{1}{|\mathcal{M}'|} \sum_{m \in \mathcal{M}'} \KL\!\left(p^{(m)}_{\text{ref}} \,\big\|\, p^{(m)}_{\text{pred}}\right),
    \label{eq:meth-kl}
\end{equation}
% SOURCE: transcribed from avg_kl_for_fill in diagnostics/run_revision.py (F.kl_div with
% reduction="batchmean" over the valid masked positions; the same quantity base_sampler logs).
which is exactly the quantity logged by the sampler in \texttt{core/base\_sampler.py} and recomputed in the revision diagnostics. By construction the ground-truth fill gives $\overline{\KL} = 0$, since then $p^{(m)}_{\text{pred}} = p^{(m)}_{\text{ref}}$ at every position, so the metric has a hard floor of zero that a perfect recovery attains. This metric is chosen over exact-match accuracy for the reason given in Section~\ref{sec:meth-task}: a recovered token that differs from the reference may still be perfectly coherent, and a metric that penalized it would confound the sampler's quality with the arbitrariness of the reference. A divergence between predictive distributions measures contextual fit directly and is insensitive to which particular admissible token was chosen. One limitation of the divergence measure must be acknowledged, because it shapes how the results are argued: its magnitude depends on the position of the corrupted token within the sequence, since corrupting an early token perturbs the conditional probability of a long suffix and yields a large divergence, whereas corrupting a final token yields almost none. For this reason no model is ranked on small differences in KL divergence; the central results are argued from the \emph{absence} of a gap between conditions rather than from a small measured gap, and a null gap is robust to this positional dependence in a way that a narrow measured win would not be. To make ``absence of a gap'' a decision rather than an impression, the analysis fixes an equivalence margin in advance, at five percent of the policy method's mean KL for each configuration, and treats two methods as practically equivalent when their paired mean difference falls inside that margin. The margin is calibrated to the pipeline's own noise rather than chosen to fit the result: as Section~\ref{sec:meth-repro} reports, the across-seed standard deviation of the final KL for the flagship configuration is $0.183$, comfortably inside the corresponding margin of about $0.327$, so the threshold below which a difference is deemed negligible is smaller than the study's stated ceiling on it yet still larger than a single standard deviation of seed noise.
% SOURCE: margin rule "0.05 x policy_mean_kl (per config)" from
% results/revision/rev_stats_gpt2.json; seed sd 0.183 vs margin 0.327 (concern 16).

Finally, for the generation-quality analyses, the study records \textbf{repetition rate}, the fraction of repeated $n$-grams, and \textbf{distinct-$n$}. These are the standard instruments for detecting the degeneration studied by \citet{holtzman2020curious}, and they are therefore the appropriate measures for the likelihood trap. For the constrained-generation extension the study records classifier accuracy and \textbf{BERTScore} \citep{zhang2020bertscore}, the latter chosen because it measures semantic similarity through contextual embeddings rather than surface overlap, and so does not penalize a valid paraphrase of the reference.

\subsection{GFlowNet Training}
\label{sec:meth-gfn}

The GFlowNet-tuned variants are produced following the amortized-inference formulation of \citet{hu2024amortizing}, using their modified sub-trajectory-balance objective \citep{malkin2022trajectory, madan2023learning}. The GFlowNet is chosen as the amortized alternative, over reward-maximizing reinforcement learning, for a reason specific to the question of this thesis: a reward-maximizing agent would converge on the single highest-reward sequence and so could not distinguish a sampler that explores the energy from one that merely finds its mode, whereas a GFlowNet is trained to sample in proportion to reward and therefore tests amortized \emph{sampling} rather than amortized optimization. It only ever evaluates the reward and never differentiates it, so the gradient behaviour that defeats the Langevin samplers is invisible to it, which is what makes it a suitable instrument for asking whether that behaviour can be sidestepped.

The policy is a low-rank adapter \citep{hu2022lora} on the frozen GPT-2 Large base, and the reward is the log-likelihood assigned by that frozen base to the completed sequence, so that the GFlowNet is trained to sample in proportion to the model's own likelihood, exactly the energy the Langevin samplers navigate. The task is a story-infilling formulation in which, given a preceding context and a known ending, the policy generates a connecting span, and the reward scores the concatenation of context, generated span, and ending under the frozen model. A single hyperparameter, denoted \texttt{len\_beta}, controls the length normalization of the reward: at $\texttt{len\_beta}=0$ the reward is the raw, unnormalized log-likelihood, and at $\texttt{len\_beta}=1$ it is fully normalized by sequence length. The consequences of this single knob, which turn out to be dramatic and to produce two opposite degeneracies, are one of the central findings of the GFlowNet study and are developed in Section~\ref{sec:results-gfn}.

Using the frozen log-likelihood as the reward may seem to contradict the thesis's own argument that this likelihood is a poor objective, so the choice is worth defending. First, that the gradient of the likelihood is uninformative does not make the likelihood value-uninformative: the gradient-free baselines of Section~\ref{sec:results-baselines} recover coherent text by evaluating exactly this quantity, so the reward carries usable signal even where its gradient does not. Second, the likelihood trap concerns the behaviour of the likelihood at its extreme, the degenerate text that globally maximizes it, not its ranking of fluent candidates near the data manifold, which is the regime the reward operates in during training. Third, and decisively, training a policy to sample in proportion to the base model's own likelihood is the standard amortized-inference formulation of \citet{hu2024amortizing}, and its premise, that this likelihood defines a target worth sampling from, is exactly the premise this thesis sets out to test; using the same reward is therefore the condition under which the test is meaningful, not an oversight.

One difference in task formulation must be made explicit, because it governs how the GFlowNet and the samplers can be compared. The GFlowNet reformulates infilling as left-to-right generation of the blank span under a restructured prompt that places the context and the known ending around it, whereas the Langevin samplers perform in-place substitution at the masked positions with bidirectional context. The two therefore solve differently shaped problems, and the comparison between them is run at the level of the energy, the final KL divergence under the same frozen base model, rather than at the level of the task. The left-to-right reformulation is not adopted for the samplers themselves, because generating the span autoregressively would dissolve the very revision capability under study, the ability to edit a token in light of both sides, and so would change the question rather than repair the metric.

\subsection{Constrained Generation Extension}
\label{sec:meth-constrained}

To test the plug-and-play claim of \eqref{eq:intro-energy} directly, which is the object of RQ4, a constrained-generation extension augments the energy with a sentiment constraint, in the manner of the gradient-based constrained samplers of \citet{kumar2022gradient} and \citet{qin2022cold}. A sentiment classification head is trained on the representations of the GPT-2 Large base and used to construct the constraint term $C(\bm{x})$. Generation is then run under five constraint modes designed to isolate the contribution of the constraint gradient. The \emph{lm\_only} mode uses only the fluency energy, providing a baseline with no steering. The \emph{cons\_only} mode uses only the constraint, showing what the constraint can do in isolation. The \emph{full} mode combines them with weights $\beta_{\text{LM}} = 0.8$ and $\beta_{\text{C}} = 0.2$, which is the intended plug-and-play configuration. The \emph{random} mode replaces the combined gradient with noise, and the \emph{cons\_random} mode replaces only the constraint gradient with noise while retaining the fluency gradient, so that the specific contribution of the constraint direction can be separated from that of the fluency direction. The steering effect is measured as the change, in percentage points, in the fraction of $n = 150$ generations per arm classified as the target sentiment relative to the unconstrained baseline. Both the discrete sampler in the formulation of this thesis and a MuCoLa-style continuous baseline are evaluated, on both infilling and continuation, so that the finding does not depend on a single sampler or task. The analysis plan is the paired point-estimate contrast \emph{cons\_only} minus \emph{cons\_random}, which cancels the shared sentiment drift; because the result files record only aggregate target-hit fractions and not per-generation labels, per-sample paired gains and their confidence intervals are not available and none are reported.

\subsection{Diagnostic Experiments}
\label{sec:meth-diagnostics}

The configurations above establish \emph{that} the methods fail. They do not, on their own, establish \emph{why}, and the why is the contribution. Five additional diagnostic experiments were therefore designed, after the main results were in hand and in direct response to them, each targeting one candidate mechanism. Three further controls were added later still, in response to the questions the first round of results raised: a sweep over step size and proposal temperature (Section~\ref{sec:results-sweep}), which establishes whether the null survives at configurations where the gradient term actually shapes the proposal; a substitution of the one-hot surrogate into the same sampler (Section~\ref{sec:results-onehot}); and a bidirectional masked language model in the same exact-energy chain (Section~\ref{sec:results-mlm}), which separates conditioning direction from training objective in the diffusion comparison. Unlike the configurations above, these do not sample; they measure properties of the energy landscape and of the gradient directly, and each is introduced here together with the question that motivated its construction.

The \textbf{linearization experiment} was built to answer the question raised by the gradient-versus-random result: if the gradient carries no usable information, is that because the first-order surrogate the discrete proposal relies on fails to predict the true change in energy? For each of $200$ sequences and a masked position, the exact gradient at that position is computed, and for a stratified set of $2000$ candidate tokens the surrogate predicted change $\hat{\Delta}(v) = \vg^\top(\ve(v) - \ve(x_i))$ is compared against the true change $\Delta(v) = \log p(\bm{x}_{i\to v}) - \log p(\bm{x})$ obtained by an exact forward pass, with the embedding distance of each candidate recorded so that the correlation can be examined as a function of distance. The true change is further decomposed into a \emph{self} term, the change in the log-probability of the candidate given its own left context, and a \emph{future} term, the change in the log-probability of the suffix, because the gradient with respect to the input embedding is structurally able to register only the latter, and measuring both quantifies how much signal the gradient is blind to.

The \textbf{acceptance experiment} was built to answer the question raised by the Metropolis--Hastings breakdown: when the correction rejects nearly every useful move on the continuous sampler, is the rejection driven by the target term, meaning the moves are genuinely bad, or by the proposal term, meaning the moves are good but the correction cannot verify their reversibility? For every proposal it logs whether the proposal crossed a Voronoi cell boundary, whether it was accepted, and the decomposition of the acceptance ratio into its target and proposal terms, which lets the cause be located unambiguously.

The \textbf{trajectory experiment} logs the raw state of the sampler at every step for a small number of sequences, so that the path can be projected into two dimensions and inspected, and it records the distance from the continuous state to the nearest token embedding, which quantifies how far the sampler drifts from the token manifold over the course of a run. The \textbf{likelihood-trap experiment} asks whether low energy is a good place to be at all. It decodes with several strategies, from greedy and beam search to ancestral sampling, and measures the relationship between the per-token likelihood of the generated text and its repetition rate. A final \textbf{anisotropy analysis} measures the distribution of inter-token distances and related geometric quantities in the two embedding spaces, underpinning both the step-size result and the unreliability of Euclidean distance.

The scale of the configuration grid, one hundred and forty-five sampling configurations each evaluated on two hundred sequences, together with the separate diagnostic runs across five models, required an experimental infrastructure that could distribute work reliably across several machines and recover cleanly from interruptions. The runs are organized as a manifest of jobs, each annotated with its memory requirement, and dispatched by a resumable file-queue system that assigns jobs to available devices according to that requirement, so that a large model and a small one are not placed on the same device when their combined memory would exceed it. Completed jobs are marked so that an interrupted sweep resumes rather than restarts, which matters when a single sweep runs for many hours. All experiments use this system, and all sampling code was verified against a reference implementation with a bitwise-equivalence test suite before the runs reported here were launched, so that the diagnostic instrumentation is known not to have altered the behaviour of the samplers it measures. One limitation of the design is worth recording, because it was encountered rather than anticipated. A job is claimed by creating a lock directory and the lock is released only on clean completion or on an explicit failure, so a worker killed abruptly, by an out-of-memory kill on its host for instance, strands its claims: the queue then appears busy while nothing runs. Recovering from this requires distinguishing a stalled job from a slow one, which the current design cannot do; a periodic heartbeat on the lock, treating a lock older than some interval as reclaimable, would remove the failure mode. The blanket reset utility is not a substitute, since it would also release jobs still running on other hosts. A second limitation is one of throughput rather than correctness: the grid runner dispatches whole configurations, so a single configuration is pinned to a single device and a large model runs at the speed of one accelerator no matter how many are idle. Splitting a configuration's sequences across devices and recombining them afterwards is straightforward, since the sample set can be resolved before any model is loaded and then partitioned, and doing so leaves per-sequence results unchanged; the diffusion experiments of Section~\ref{sec:results-diffusion} were run that way, and the autoregressive grid was not. The Llama-3 and GFlowNet families would have completed several times faster had they been.

\subsection{Reproducibility, Seeds, and Hardware}
\label{sec:meth-repro}

Because the central claims of this thesis are claims about the \emph{absence} of an effect, the run-to-run stability of the pipeline is itself part of the evidence and is quantified directly rather than assumed. Rerunning the flagship configuration, the discrete sampler on GPT-2 Large with the Metropolis--Hastings correction and gradient normalization enabled over a fifty-step schedule, under four independent corruption seeds, $1000$ through $1003$, yields final KL divergences of $6.348$, $6.589$, $6.150$, and $6.291$, with a mean of $6.345$, a standard deviation of $0.183$, and a range of $0.438$.
% SOURCE: REVISION_LOG.md concern 16, seeds 1000-1003 of
% gpt2-large.dls.policy.mh.gn.free.s50: mean 6.345, sd 0.183, range 0.438.
This across-seed standard deviation fixes the scale below which a difference in KL cannot be told apart from noise, and it is the empirical basis both for the equivalence margin defined in Section~\ref{sec:meth-metrics} and for the bounded-effect reading of the gradient-versus-random comparisons in Section~\ref{sec:results-fallacy}. The corruption procedure is deterministic per sample index, with the seed set to the data seed plus the running index, so the corrupted inputs, and therefore the paired comparisons, are identical across the three method runs of a configuration and can be reproduced bit for bit; this is what allows the statistical analysis of Section~\ref{sec:results-fallacy} to pair the policy method against its comparators on the same sample without any rerun.

All reported experiments were run on a university GPU server equipped with $49$\,GB A6000-class accelerators, under PyTorch~$2.12$ with CUDA~$13.0$. For code and data availability, the sampler and diagnostic implementations, together with the bitwise-equivalence suite (\texttt{verify\_equivalence\_suite.py}) that certifies the instrumented code against the reference, and the manifests that define every configuration in the grid, are retained with the thesis so that the full set of runs can be regenerated from the seeds above. The accompanying \texttt{README.md} documents the environment, the exact rerun command for every experiment family, and an artifact map that links each table and figure in this thesis to the script that produces it and the result file it draws from, so that every quantitative claim is traceable from the number in the text to the file on disk.
